Drug-target interaction (DTI) papers often use synthetic cold-start splits as proxies for generalization, but it remains unclear whether conclusions drawn from these splits transfer to real distribution shifts. We present a benchmark-first study showing that model rankings are highly sensitive to how the shift is defined. Using a compact panel consisting of a plain base encoder, a retrieval-augmented model (\texttt{RAICD}), and a functional target-memory variant (\texttt{FTM sparse}), we first build a 3-seed synthetic reference panel on \texttt{BindingDB\_Kd}. Even within this single dataset, rankings are heterogeneous: \texttt{RAICD} improves over the base model on \texttt{blind\_start}, loses on \texttt{unseen\_drug}, and \texttt{FTM sparse} is strongest on synthetic \texttt{unseen\_target}. We then evaluate the same panel on a real temporal/provenance benchmark, \texttt{BindingDB\_patent / patent\_temporal}, where the internal-panel ranking reverses. Over 3 seeds, the plain base model achieves $0.7772 \pm 0.0086$ \AUPRC{}, outperforming both \texttt{RAICD} ($0.7692 \pm 0.0018$) and \texttt{FTM sparse} ($0.7721 \pm 0.0025$). When we expand the model panel with recent pooled-LM baselines, \texttt{DTI-LM} becomes the strongest patent model at $0.7825 \pm 0.0078$ \AUPRC{}, and the same winner is preserved under an earlier alternative patent cutoff. Paired bootstrap intervals confirm the sign of the key \AUPRC{} deltas. A year-band and overlap-bucket diagnosis further shows that temporal OOD is structured rather than monolithic. Matched 3-seed support from \texttt{BindingDB\_Ki} and \texttt{DAVIS} indicates that synthetic \texttt{unseen\_target} gains are not reliably portable, and we package the split manifests and evaluator as reusable benchmark resources. Within the studied model panel and benchmarks, these results show that synthetic cold splits and real OOD shifts are not interchangeable evaluation settings for DTI.
